La siguiente tesina se organiza de la siguiente manera:

\begin{itemize}
    \item \textbf{Capítulo 2:} Desarrolla el marco teórico y el estado del arte. Aquí se introducen los conceptos matemáticos y computacionales fundamentales, así como las herramientas y antecedentes necesarios para comprender el proyecto.
    
    \item \textbf{Capítulo 3:} Detalla la implementación técnica del simulador. Se documenta el desarrollo del modelo físico-matemático de la dinámica vehicular, la construcción del entorno de visualización y la programación de la biblioteca de interfaces para la integración con el hardware.
    
    \item \textbf{Capítulo 4:} Presenta la validación del sistema y los resultados. Se exponen casos de estudio con el simulador en pleno funcionamiento, evaluando su rendimiento y estabilidad tanto en simulaciones precalculadas (lazo abierto) como en ejecuciones interactivas en tiempo real (lazo cerrado).
    
    \item \textbf{Capítulo 5:} Expone las conclusiones finales del trabajo. Se analizan los resultados obtenidos en contraste con los objetivos planteados, se discuten las limitaciones del modelo desarrollado y se proponen posibles líneas de trabajo futuro para la extensión y mejora de esta herramienta.
\end{itemize}